(1) 均匀各向同性良导体占 \(x>0\)，电导率 \(\sigma>0\)，磁导率 \(\mu>0\)，介电常数 \(\varepsilon\)，在角频率 \(\omega>0\) 下满足 \(\sigma\gg\omega\varepsilon\)。表面给定切向谐波磁场 \(\mathbf H(0,t)=\operatorname{Re}[H_0e^{-i\omega t}]\hat z\)，深处场消失。忽略位移电流，由Maxwell方程推导导体内磁扩散方程与衰减场 \(H_z(x)\)。
(2) 求复电流密度 \(j_y(x)\)、电场 \(E_y(x)\)，定义表面阻抗 \(Z_s=E_y(0)/H_z(0)\)，给出实部、虚部和采用 \(e^{-i\omega t}\) 约定的相位符号。
(3) 用复Poynting矢量求穿过表面的时间平均输入功率密度，再独立积分全半空间的体焦耳热，证明两者严格相等。
(4) 讨论 \(\omega\to0^+\)、\(\omega\to\infty\) 时趋肤深度、阻抗与功率的标度，说明高频近似何时因位移电流、非局域输运或表面粗糙失效。
